\section{A kernel-clustering application}\label{kc:section}

\subsection{The objective and the approximation threshold}

This application uses Theorem~\ref{thm:main},
the kernel-clustering results of Khot and Naor, and the companion paper
\emph{The Unique Games Theorem}. None of these complexity arguments is used
in the proof of the Gaussian inequality.

Fix an integer \(k\ge3\). An input is an explicitly represented rational
symmetric matrix \(A=(a_{pq})_{p,q=1}^N\) satisfying
\(A\succeq0\) and \(A\one=0\). The latter is the centering condition;
for a positive semidefinite matrix it is equivalent to
\(\sum_{p,q}a_{pq}=0\). For a map \(\sigma:[N]\to[k]\), define
\begin{equation}\label{kc:value}
 \val_A(\sigma)=\sum_{p,q=1}^N a_{pq}\one_{\{\sigma(p)=\sigma(q)\}},
 \qquad
 \Clust(A\mid I_k)=\max_{\sigma:[N]\to[k]}\val_A(\sigma).
\end{equation}
Here \([m]=\{1,\ldots,m\}\), and \(I_k\) is the \(k\)-by-\(k\)
identity target matrix. The sum is over ordered pairs and includes the
diagonal. Equivalently, it is
\(\sum_{i=1}^k\sum_{p,q\in S_i}a_{pq}\) for the cells
\(S_i=\sigma^{-1}(i)\); empty cells are allowed. There is no
normalization by cell cardinalities. Each objective value is nonnegative,
since it is a sum of quadratic forms \(\one_{S_i}^{\mathsf T}A\one_{S_i}\).

A deterministic approximation with \emph{loss factor} \(\alpha\ge1\)
returns an assignment of value at least
\(\Clust(A\mid I_k)/\alpha\). An expected approximation has the same
inequality with the returned value replaced by its expectation.
Put
\[
 c_*=\frac9{8\pi},
 \qquad
 \alpha_k=\frac{1-1/k}{c_*}
          =\frac{8\pi}{9}\left(1-\frac1k\right).
\]

\paragraph{The matching rounding constant.}
Khot and Naor's approximation theorem
\cite[Theorem~2.1]{KhotNaor}
gives the expected rounding bound
\[
 \E\val_A(\sigma)\ge\frac{\Clust(A\mid I_k)}{\alpha_k}
\]
on the inputs in \eqref{kc:value}. This is their endpoint analytic
guarantee using an optimal semidefinite relaxation and ideal Gaussian
sampling. We quote this rounding guarantee separately from the
rational-input hardness statement below; an exact finite-bit
implementation at the endpoint is not asserted here.

\begin{theorem}[Kernel-clustering hardness]\label{kc:main}
For each fixed \(k\ge3\) and every fixed
\(1\le\alpha<\alpha_k\), it is NP-hard, given a rational symmetric
matrix \(A\succeq0\) with \(A\one=0\) and the identity-target
objective \eqref{kc:value}, to find
\(\sigma:[N]\to[k]\) such that
\[
 \val_A(\sigma)\ge\Clust(A\mid I_k)/\alpha.
\]
In particular, such a deterministic polynomial-time approximation
would imply \(\mathrm P=\mathrm{NP}\).
\end{theorem}

We prove the hardness statement below. It concerns strict improvement
over \(\alpha_k\), not hardness at the endpoint.

For the more general objective $\sum_{p,q}a_{pq}b_{\sigma(p)\sigma(q)}$
with a fixed positive semidefinite target matrix $B=(b_{ij})_{i,j=1}^k$,
Khot and Naor later characterized the Unique Games hardness threshold
for approximating the optimal clustering value through geometric and
Gaussian partition parameters
\cite[Theorem~1.1]{KhotNaorSharp}. Here the target is $I_k$, and the
Gaussian constant is supplied by Theorem~\ref{thm:main}.

\subsection{The two external hardness inputs}

The Gaussian parameter in Khot and Naor's theorem is
\[
 C_k:=\sup_{(B_1,\ldots,B_k)}
 \sum_{i=1}^k
 \left\|\int_{B_i}x\,d\gamma_{k-1}(x)\right\|^2.
\]
The supremum ranges over measurable partitions of \(\R^{k-1}\),
with empty cells allowed. Since \(k\ge3\),
Theorem~\ref{thm:main} and the attaining example in
Section~\ref{sec:example} give
\begin{equation}\label{kc:gaussian-constant}
 C_k=c_*.
\end{equation}

We first specify the discrete analytic input. Let \(\Omega=[k]\) have
uniform probability measure. Choose a real orthonormal basis
\(b_0,b_1,\ldots,b_{k-1}\) of functions on \(\Omega\), with \(b_0=1\).
For \(r\in\N\) and \(\nu\in\{0,\ldots,k-1\}^r\), write
\[
 b_\nu(x)=\prod_{j=1}^r b_{\nu_j}(x_j),
 \qquad
 |\nu|=\bigl|\{j:\nu_j\ne0\}\bigr|.
\]
Every \(h:\Omega^r\to\R\) has expansion
\(h=\sum_\nu\widehat h(\nu)b_\nu\), with coefficients computed using
uniform probability measure. For real \(D\ge0\), set
\begin{equation}\label{kc:influence}
 \Inf_j^{\le D}(h)=
 \sum_{\substack{|\nu|\le D\\\nu_j\ne0}}\widehat h(\nu)^2.
\end{equation}
Thus a nonintegral cutoff is interpreted by the displayed inequality
on the integer \(|\nu|\). Let \(P_1\) be orthogonal projection onto
the span of \(b_\nu\) with \(|\nu|=1\), and, for
\(f=(f_1,\ldots,f_k)\), put
\[
 \OBJ(f)=\sum_{i=1}^k\|P_1f_i\|_2^2.
\]
All \(L_2\) norms in this section use uniform probability measure.
Define the closed simplex with slack by
\[
 \Delta_k=\{u\in\R^k:u_i\ge0,\ \sum_{i=1}^k u_i\le1\}.
\]

\begin{theorem}[Khot--Naor low-influence bound]\label{kc:kn-input}
For fixed \(k\) and every \(\eta>0\), there exists \(0<\tau_0<1/2\),
independent of \(r\), such that every \(r\in\N\) and every
\(f:\Omega^r\to\Delta_k\) satisfying
\[
 \Inf_j^{\le\log(1/\tau_0)}(f_i)\le\tau_0
 \quad(i\in[k],\ j\in[r])
\]
obey \(\OBJ(f)\le C_k+\eta\).
\end{theorem}

This is \cite[Theorem~3.8]{KhotNaor}, with the same uniform product
measure, Fourier degree, simplex, and Gaussian normalization. Its proof
uses the invariance principle of Mossel, O'Donnell, and
Oleszkiewicz~\cite{MOO} to compare low-influence functions on product
spaces with their Gaussian counterparts.
The restriction \(\tau_0<1/2\) is harmless: decreasing the threshold
and increasing the degree cutoff strengthen the hypotheses.
Choose a rational \(0<\tau<\min\{\tau_0,1/2\}\) and set
\begin{equation}\label{kc:cutoff}
 D=\max\{1,\lceil\log(1/\tau)\rceil\}.
\end{equation}
Monotonicity of \eqref{kc:influence} in the cutoff gives the useful
contrapositive
\begin{equation}\label{kc:contrapositive}
 \OBJ(f)>C_k+\eta
 \quad\Longrightarrow\quad
 \Inf_j^{\le D}(f_i)>\tau
 \text{ for some }i\in[k],\ j\in[r].
\end{equation}
Indeed, if every influence at cutoff \(D\) were at most \(\tau\),
every influence at the original cutoff would be at most \(\tau_0\).

The second input is the following exact interface of
\cite[Theorem~1.1]{UniqueGamesTheorem}.

\begin{theorem}[Unique Games input]\label{kc:ug-input}
For every fixed \(\varepsilon,\delta\in(0,1/2)\), there are an
integer \(s\ge1\) and a deterministic polynomial-time reduction
from \(\mathrm{3SAT}\) to explicit unweighted simple bipartite
Unique Games instances over \(\mathbb F_2^s\). If the formula is
satisfiable the game has a labeling satisfying at least
\(1-\varepsilon\) of its edges; if it is unsatisfiable every
labeling satisfies at most \(\delta\) of its edges. The edge set
is nonempty and each constraint is a translation of the alphabet.
The alphabet size and the polynomial's degree and constants depend
only on \(\varepsilon,\delta\).
\end{theorem}

Only ordinary satisfied-edge fractions are asserted here. We will
use neither regularity nor the stronger condition that almost every
left vertex has all its incident constraints satisfied.

\subsection{A rational centered positive semidefinite matrix}

Let \(G=(V,W,E)\) be a game supplied by
Theorem~\ref{kc:ug-input}, and enumerate its alphabet as
\([r]\), where \(r=2^s\). Reorient constraints if necessary so that
the permutation \(\pi_{vw}:[r]\to[r]\) attached to \(vw\in E\)
is satisfied precisely when
\[
 \ell(v)=\pi_{vw}(\ell(w)).
\]
Translations are permutations, so this changes no game value.
Delete isolated vertices, put \(M=|E|\), and write
\[
 d_v=|N(v)|,\qquad \mu(v)=\frac{d_v}{M}\quad(v\in V).
\]
Choosing \(v\) with probability \(\mu(v)\), then a uniform
\(w\in N(v)\), chooses an edge uniformly.

To fix the permutation convention, for a permutation \(\pi\) define
\[
 (R_\pi x)_a=x_{\pi(a)}\quad(a\in[r]).
\]
For a scalar table \(h:W\times\Omega^r\to\R\), define
\begin{equation}\label{kc:averaging}
 (T_vh)(x)=\frac1{d_v}
       \sum_{w\in N(v)}h(w,R_{\pi_{vw}}x),
 \qquad
 Q(h)=\sum_{v\in V}\mu(v)\|P_1T_vh\|_2^2 .
\end{equation}
The output matrix will represent \(Q\) in the ordinary,
unnormalized Euclidean coordinates of the table \(h\).

\begin{lemma}\label{kc:matrix}
For fixed \(k,r\), one can construct in deterministic polynomial
time a rational symmetric matrix \(A\), indexed by
\(W\times\Omega^r\), such that
\[
 h^{\mathsf T}Ah=Q(h),\qquad A\succeq0,\qquad A\one=0 .
\]
For every \(\sigma:W\times\Omega^r\to[k]\), if
\(F_i(w,x)=\one_{\{\sigma(w,x)=i\}}\), then
\begin{equation}\label{kc:objective-identity}
 \val_A(\sigma)=
 \sum_{v\in V}\mu(v)\OBJ(F_v),\qquad
 F_v=(T_vF_1,\ldots,T_vF_k):\Omega^r\to\Delta_k.
\end{equation}
All output assignment values lie in \([0,1]\).
\end{lemma}

\begin{proof}
Put \(q=k^r\). The projection \(P_1\), written as a matrix acting
on the counting-coordinate vector \((h(x))_{x\in\Omega^r}\), is
\begin{equation}\label{kc:projection-matrix}
 (P_1)_{x,y}=\frac1q
       \sum_{j=1}^r\bigl(k\one_{\{x_j=y_j\}}-1\bigr).
\end{equation}
To see this, the centered subspaces of functions of distinct
individual coordinates are mutually orthogonal. Hence
\[
 (P_1h)(x)=
 \sum_{j=1}^r\bigl(\E[h(X)\mid X_j=x_j]-\E h(X)\bigr),
\]
which is \eqref{kc:projection-matrix}.
This matrix is rational, symmetric and idempotent.
The counting-coordinate matrix of \(T_v\) has entries
\[
 (T_v)_{x,(w,y)}
 =\frac1{d_v}\one_{\{w\in N(v)\}}
           \one_{\{y=R_{\pi_{vw}}x\}} .
\]
Define
\begin{equation}\label{kc:matrix-formula}
 A=\frac1q\sum_{v\in V}\mu(v)T_v^{\mathsf T}P_1T_v.
\end{equation}
There are two factors \(q^{-1}\): one in
\eqref{kc:projection-matrix}, and one outside the sum in
\eqref{kc:matrix-formula} because \(L_2\) uses probability
measure. Symmetry and \(P_1^2=P_1\) give
\[
 h^{\mathsf T}Ah
 =\sum_v\mu(v)\frac1q\|P_1T_vh\|_{\ell_2}^2=Q(h)\ge0.
\]
Also \(T_v\one=\one\) and \(P_1\one=0\), so
\eqref{kc:matrix-formula} gives \(A\one=0\).

There are \(N=|W|k^r\) output indices. For fixed \(k,r\) this is
polynomial in the explicit game size, and the dense matrix can be
computed by rational arithmetic in polynomial time. Its entries
are sums of terms with denominators dividing
\(M d_v q^2\). A common denominator divides
\(M q^2\prod_{v\in V}d_v\), whose bit length is
\(O(\log M+\sum_v\log d_v+\log q)\), and the numerators also
have polynomial bit length. No irrational Fourier basis needs
to be constructed.

By definition of the ordered sum,
\(\val_A(\sigma)=\sum_iF_i^{\mathsf T}AF_i\); inserting
\eqref{kc:averaging} proves \eqref{kc:objective-identity}.
Every \(F_v(x)\) is an average of coordinate unit vectors, so
its entries are nonnegative and sum to one. Parseval and
orthogonal projection give
\begin{equation}\label{kc:objective-bound}
 0\le\OBJ(F_v)
 \le\sum_{i=1}^k\|F_{v,i}\|_2^2
 \le\E\sum_{i=1}^k F_{v,i}=1.
\end{equation}
Since \(\sum_v\mu(v)=1\), the last claim follows.
\end{proof}

\subsection{Completeness and weighted decoding}

The next two lemmas compare ordinary edge-value completeness and
soundness with the kernel objective. They adapt the
Khot--Naor reduction \cite[Section~3.3]{KhotNaor} to the precise
input of Theorem~\ref{kc:ug-input}.

\begin{lemma}[Completeness]\label{kc:completeness}
If a labeling of \(G\) satisfies at least \(1-\varepsilon\) of
its edges, then the matrix of Lemma~\ref{kc:matrix} satisfies
\[
 \Clust(A\mid I_k)\ge
       \left(1-\frac1k\right)(1-\varepsilon)^2.
\]
\end{lemma}
\begin{proof}
Use a witnessing labeling \(\ell\) to set
\(\sigma(w,x)=x_{\ell(w)}\).
For \(j\in[r]\), let \(D_j(x)=e_{x_j}\), where
\(e_1,\ldots,e_k\) are the coordinate unit vectors in \(\R^k\).
With
\[
 p_{v,j}=
 \Pr_{w\in N(v)}\{\pi_{vw}(\ell(w))=j\},
\]
the chosen permutation convention gives
\(F_v=\sum_{j=1}^r p_{v,j}D_j\).
For each \(i\), the level-one part of the \(i\)-th component of
\(D_j\) is \(\one_{\{x_j=i\}}-1/k\). These parts for different
\(j\)'s are orthogonal. Moreover,
\[
 \sum_{i=1}^k
 \E\bigl(\one_{\{X_j=i\}}-1/k\bigr)^2=1-\frac1k.
\]
It follows that
\[
 \OBJ(F_v)=
 \left(1-\frac1k\right)\sum_{j=1}^r p_{v,j}^2
 \ge \left(1-\frac1k\right)p_{v,\ell(v)}^2.
\]
The weighted mean \(\E_{v\sim\mu}p_{v,\ell(v)}\) is exactly
the satisfied-edge fraction of \(\ell\), and is therefore at
least \(1-\varepsilon\). Jensen's inequality and
\eqref{kc:objective-identity} prove the claim.
\end{proof}

\begin{lemma}[Soundness]\label{kc:soundness}
Fix \(\eta>0\) with \(C_k+2\eta<1\), and choose
\(\tau,D\) as in \eqref{kc:cutoff} for this \(\eta\).
If every labeling of \(G\) satisfies at most \(\delta\) of its
edges, where
\begin{equation}\label{kc:delta}
 0<\delta<\frac{\eta\tau^2}{4kD},
\end{equation}
then
\(\Clust(A\mid I_k)\le C_k+2\eta\).
\end{lemma}
\begin{proof}
Suppose an assignment \(\sigma\) gives a larger value, and use
the functions in Lemma~\ref{kc:matrix}.
Let
\[
 V_{\mathrm{good}}=
       \{v:\OBJ(F_v)>C_k+\eta\}.
\]
If its \(\mu\)-mass is \(\beta\), then
\eqref{kc:objective-bound} gives
\[
 \sum_v\mu(v)\OBJ(F_v)
 \le\beta+(1-\beta)(C_k+\eta)
 \le C_k+\eta+\beta.
\]
The assumed objective therefore forces \(\beta>\eta\).

For every \(v\in V_{\mathrm{good}}\),
\eqref{kc:contrapositive} supplies
\(i_v\in[k]\) and \(j_v\in[r]\) with
\(\Inf_{j_v}^{\le D}(F_{v,i_v})>\tau\).
Taking the squared norm of the relevant Fourier coefficients
in the neighbor average, and using convexity, gives
\begin{equation}\label{kc:influence-average}
 \Inf_{j_v}^{\le D}(F_{v,i_v})
 \le \frac1{d_v}\sum_{w\in N(v)}
       \Inf_{\pi_{vw}^{-1}(j_v)}^{\le D}(F_{i_v}(w,\cdot)).
\end{equation}
The inverse permutation in this formula follows directly from
\((R_\pi x)_a=x_{\pi(a)}\): the Fourier indices involving
left coordinate \(j_v\) come from right coordinate
\(\pi^{-1}(j_v)\).
Every influence on the right belongs to \([0,1]\), by
Parseval and \(0\le F_i(w,\cdot)\le1\).
An average exceeding \(\tau\) has more than a \(\tau/2\)
fraction of its entries at least \(\tau/2\):
otherwise its average is at most
\(\tau/2+(\tau/2)\cdot1\le\tau\).

For each \(w\in W\), form the finite label list
\[
 L(w)=
 \{a\in[r]:\Inf_a^{\le D}(F_i(w,\cdot))\ge\tau/2
                      \text{ for some }i\in[k]\}.
\]
For every scalar \(h\), counting how often a Fourier coefficient
contributes gives
\[
 \sum_{a=1}^r\Inf_a^{\le D}(h)
 =\sum_{|\nu|\le D}|\nu|\,\widehat h(\nu)^2
 \le D\|h\|_2^2.
\]
Consequently
\begin{equation}\label{kc:list-size}
 |L(w)|\frac{\tau}{2}
 \le\sum_{i=1}^k\sum_{a=1}^r
            \Inf_a^{\le D}(F_i(w,\cdot))
 \le kD,
 \qquad |L(w)|\le\frac{2kD}{\tau}.
\end{equation}

Label a good \(v\) by \(j_v\), and other left vertices arbitrarily.
Independently label each right vertex uniformly from \(L(w)\)
when this list is nonempty, and arbitrarily otherwise.
For a good \(v\), more than a \(\tau/2\) fraction of its
neighbors have \(\pi_{vw}^{-1}(j_v)\in L(w)\), by
\eqref{kc:influence-average}. On each such edge the random label
at \(w\) satisfies the constraint with probability at least
\(\tau/(2kD)\), by \eqref{kc:list-size}.
The expected fraction of satisfied edges is therefore greater than
\[
 \eta\cdot\frac{\tau}{2}\cdot\frac{\tau}{2kD}
   =\frac{\eta\tau^2}{4kD}.
\]
Here the weights \(\mu(v)=d_v/M\) ensure that this is the
uniform-edge fraction. Some deterministic labeling attains at
least its expectation, contradicting \eqref{kc:delta}.
\end{proof}

\subsection{Completing the reduction}

\begin{proof}[Proof of the hardness assertion in Theorem~\ref{kc:main}]
Fix \(k\ge3\) and \(1\le\alpha<\alpha_k\).
First choose rational \(\varepsilon,\eta>0\) sufficiently small
that \(\varepsilon<1/2\), \(c_*+2\eta<1\), and
\begin{equation}\label{kc:separated-gap}
 \alpha<
 \frac{(1-1/k)(1-\varepsilon)^2}{c_*+2\eta}.
\end{equation}
Next choose \(\tau,D\) from Theorem~\ref{kc:kn-input} and
\eqref{kc:cutoff}; these depend on \(k,\eta\), not the game size.
Choose a rational
\[
 0<\delta<
 \min\left\{\frac12,\frac{\eta\tau^2}{4kD}\right\}.
\]
Apply Theorem~\ref{kc:ug-input} with these fixed errors.
Only at this stage is the fixed alphabet size \(r=2^s\)
determined. In particular \(k^r\) is a constant with respect to
the input formula size, although no uniform bound polynomial in
the error reciprocals is asserted.

Compose that deterministic reduction with Lemma~\ref{kc:matrix}.
Lemmas~\ref{kc:completeness} and \ref{kc:soundness}, together with
\eqref{kc:gaussian-constant}, give the gap
\[
 \begin{aligned}
 \varphi\text{ satisfiable}
 &\ \Longrightarrow\
 \Clust(A_\varphi\mid I_k)\ge
       H:=(1-1/k)(1-\varepsilon)^2,\\
 \varphi\text{ unsatisfiable}
 &\ \Longrightarrow\
 \Clust(A_\varphi\mid I_k)\le
       S:=c_*+2\eta .
 \end{aligned}
\]
The output is a rational symmetric centered PSD matrix with
polynomial dimension and bit length. By \eqref{kc:separated-gap},
\(S<H/\alpha\). Choose a fixed rational
\(\theta\) strictly between these numbers.
A deterministic loss-factor-\(\alpha\) algorithm produces value
greater than \(\theta\) on satisfiable formulas; every assignment
has value less than \(\theta\) on unsatisfiable formulas.
The returned value is an exactly computable rational number.
This would decide \(\mathrm{3SAT}\) in polynomial time.
Thus achieving the stated factor is NP-hard.
\end{proof}

\begin{remark}[Randomized guarantees]\label{kc:randomized}
The quoted Khot--Naor rounding guarantee is in expectation.
The deterministic hardness conclusion of Theorem~\ref{kc:main} does
not use \(\mathrm P\ne\mathrm{NP}\) to exclude randomized
algorithms. More precisely, a polynomial-time randomized algorithm
whose expected value is at least \(\Clust(A\mid I_k)/\alpha\)
for some fixed \(\alpha<\alpha_k\) would place
\(\mathrm{3SAT}\) in \(\mathrm{RP}\).
Indeed, on the reduction instances its returned value \(X\) lies
in \([0,1]\). In the satisfiable case \(\E X\ge H/\alpha>\theta\);
therefore
\[
 \Pr\{X>\theta\}\ge
 \frac{H/\alpha-\theta}{1-\theta}>0.
\]
This is a fixed positive constant. In the unsatisfiable case
\(X\le S<\theta\) always. A fixed number of repetitions gives
one-sided success probability at least \(1/2\).
\end{remark}

\begin{remark}[Inputs to the application]\label{kc:dependencies}
The Gaussian proof supplies only \eqref{kc:gaussian-constant}.
Khot and Naor supply the expected rounding guarantee
\cite[Theorem~2.1]{KhotNaor} and the analytic
low-influence bound \cite[Theorem~3.8]{KhotNaor}.
Their reduction in Section~3.3 is the model for
Lemmas~\ref{kc:matrix}--\ref{kc:soundness}; the degree weights and
the squared completeness loss explicitly accommodate the ordinary
edge-value, possibly nonregular instances from
\cite[Theorem~1.1]{UniqueGamesTheorem}.
No stronger Unique Games promise is assumed.
The separate Unique Games theorem is a formal dependency only
of this hardness application. The algorithmic endpoint is quoted
in its ideal rounding formulation; this section supplies no
finite-precision endpoint implementation.
\end{remark}

